\section*{Overview}

Rerooting DP is the pattern for ``solve the same tree DP for every root.'' Instead of recomputing from scratch \(n\)
times, it uses one pass to gather subtree information and a second pass to transfer the parent-side contribution into
each child.

\section*{When to Use It}

Use rerooting when:

\begin{itemize}
  \item the tree answer is wanted at every vertex,
  \item the answer for one root can be updated locally when moving the root across one edge,
  \item the DP contribution can be merged from children plus one parent-side term.
\end{itemize}

\section*{Core Idea}

The first DFS computes the usual subtree DP. The second DFS tells each child what the world outside that child's
subtree contributes, so the child can combine:

\begin{itemize}
  \item its own subtree data,
  \item the contribution coming from the parent side.
\end{itemize}

\section*{Key Insight}

Moving the root from a parent to a child only changes one edge's direction. If the DP can be updated across that local
change, all-roots answers become a transfer problem rather than \(n\) separate DPs.

\section*{Operations / Main Technique}

\begin{itemize}
  \item compute downward subtree values,
  \item derive the full answer at the initial root,
  \item pass adjusted parent-side information into each child,
  \item record the answer for every node.
\end{itemize}

\paragraph{Worked example.}
For the sum of distances from each node to all others, moving the root from \(v\) to child \(to\) decreases the
distance to every node in \(to\)'s subtree by \(1\) and increases the distance to every other node by \(1\).

\section*{Correctness Intuition}

The subtree pass is ordinary tree DP. The reroot pass is correct because each child receives exactly the contribution of
everything not in its own subtree, adjusted as if that child were the new root. Together, those two pieces cover the
entire tree without overlap.

\section*{Complexity Analysis}

Most rerooting solutions run in \(O(n)\) time once the merge and transfer formulas are \(O(1)\) per edge.

\section*{Implementation}

The sample code solves the classic ``sum of distances to all nodes'' problem. It is a good template because the transfer
formula is explicit and easy to check.

\section*{Common Pitfalls}

\begin{itemize}
  \item Forgetting to remove a child's own contribution before passing the parent-side value back into it.
  \item Designing a state that works for one root but is too weak to transfer cleanly.
  \item Recomputing siblings from scratch instead of deriving the transferred value algebraically.
\end{itemize}

\section*{Variants / Extensions}

\begin{itemize}
  \item Max/min answers for every root.
  \item Prefix-suffix child merges when each child needs the aggregate of all other children.
  \item Generic rerooting templates with custom combine and lift operations.
\end{itemize}

\section*{Practice Problems}

\begin{itemize}
  \item Sum of distances from every node.
  \item Best root under a subtree-based score.
  \item Count or optimize rooted configurations for every possible root.
\end{itemize}

\section*{References}

\begin{itemize}
  \item \href{https://usaco.guide/gold/all-roots?lang=cpp}{USACO Guide: DP on Trees - Solving for All Roots}.
  \item \href{https://codeforces.com/catalog?locale=en}{Codeforces Catalog}.
  \item \href{https://github.com/kth-competitive-programming/kactl}{KACTL}.
\end{itemize}
